\begin{phasebox}
\textbf{هدف:} مشخص‌کردن جای هر فایل بدون ایجاد پوشه‌های زیاد.\\
\textbf{پیش‌نیاز:} محدوده نسخه اولیه مشخص باشد.\\
\textbf{معیار پایان:} برای هر نوع کد فقط یک پوشه واضح وجود داشته باشد.
\end{phasebox}

\section{ساختار ساده پروژه}
کد پروژه فقط در چند پوشه اصلی قرار می‌گیرد:

\begin{itemize}
    \item پوشه
    \lr{\texttt{arduino}}
    برای کد برد و سنسورها؛
    \item پوشه
    \lr{\texttt{backend}}
    برای منطق اصلی برنامه و ارتباط
    \lr{MQTT}
    در پایتون؛
    \item پوشه
    \lr{\texttt{database}}
    برای ساخت جدول‌ها و تمام کدهای دسترسی به پایگاه داده؛
    \item پوشه
    \lr{\texttt{webapp}}
    برای برنامه
    \lr{Flask}
    و صفحه‌های وب؛
    \item پوشه
    \lr{\texttt{contracts}}
    برای نام موضوع‌های پیام و قالب داده‌ها؛
    \item پوشه
    \lr{\texttt{tests}}
    برای تست ارتباط بین قسمت‌ها.
\end{itemize}

\section{ارتباط قسمت‌ها}
برد پیام را منتشر می‌کند. بک‌اند پیام را دریافت و بررسی می‌کند. کدهای پوشه پایگاه داده آن را ذخیره می‌کنند. برنامه وب برای نمایش یا کنترل، توابع بک‌اند را صدا می‌زند.

\section{قواعد ساده}
\begin{itemize}
    \item دستورهای پایگاه داده فقط در پوشه مربوط به پایگاه داده نوشته شوند.
    \item کد صفحه‌ها و مسیرهای وب فقط در پوشه برنامه وب باشند.
    \item کد سخت‌افزار فقط در پوشه آردوینو باشد.
    \item نام پیام‌ها در یک محل مشترک تعریف شوند تا پایتون و آردوینو با هم سازگار بمانند.
    \item تا زمانی که تعداد فایل‌ها زیاد نشده، پوشه داخلی جدید ساخته نشود.
\end{itemize}

\section{تغییر ابزارها در آینده}
برای جایگزینی
\lr{Flask}
با
\lr{Django}
فقط پوشه برنامه وب تغییر می‌کند. برای جایگزینی
\lr{SQLite}
با پایگاه داده دیگر، تغییر اصلی در پوشه پایگاه داده انجام می‌شود.
